% Per-game replacements of the fixed texts of lang-en.tex (\renewcommand).
\renewcommand{\tComputerText}{\textbf{An MSX1, MSX2 or turbo R} with two cartridge slots, one for the game and one for the V9990. On a turbo R the game runs at the same speed as on an MSX1.}
\renewcommand{\tMemoryText}{\textbf{\ramNeeded{} of RAM}. An MSX with less memory can still play if the flash cartridge that holds the game has its own memory (for example a MegaFlashROM SCC+ SD): the game uses it.}
\renewcommand{\tSoundText}{The MSX's own sound: title music, police siren, engine, turbo and effects. No sound cartridge is needed.}
\renewcommand{\tGameText}{The \textbf{\gameTitle{} V9990 cartridge}: a \romKind{} (the file \texttt{\romFile}, written to a flash cartridge or equivalent). It works with the cartridge in a normal slot or in a sub-slot.}
\renewcommand{\tControlsText}{The cursor keys and the keyboard, or a joystick with two buttons in port 1.}
\renewcommand{\tSetJoy}{Plug a joystick into port 1 if you want to use one. Button 2 fires the turbo.}
\renewcommand{\tSetOn}{Switch the MSX on. For about two seconds the MSX's own screen shows \texttt{OCEAN SOFTWARE 1989}, \texttt{DIHALT STUDIO 2026} and \texttt{V9990 found !}; then the picture appears on the V9990's monitor. If it says \texttt{V9990 NOT found!!!} the V9990 is missing or not seated, and the game stops there.}
% Credits come after the game sections 5-8.
\renewcommand{\tCreditsTitle}{9. Credits and rights}
\newcommand{\tCreditsExtra}{%
\item The original \textbf{Chase H.Q.} is an arcade game by \textbf{Taito Corporation}, 1988. The MSX version by Ocean: program Jobbee, graphics Bill, music Jon Dunn.
\item The cover illustration is the artwork of the Taito arcade game. The music and the siren follow the ZX Spectrum 128K version.}
